\documentclass[11pt,a4paper]{article}
\usepackage[T1]{fontenc}
\usepackage{lmodern,amsmath,amssymb,amsthm,mathtools}
\usepackage[margin=27mm]{geometry}
\usepackage[colorlinks=true,linkcolor=blue,urlcolor=blue]{hyperref}
\newtheorem{theorem}{Theorem}
\newtheorem{lemma}[theorem]{Lemma}
\newtheorem{corollary}[theorem]{Corollary}
\newtheorem{proposition}[theorem]{Proposition}
\newcommand{\Q}{\mathbb Q}\newcommand{\Z}{\mathbb Z}
\newcommand{\ord}{\operatorname{ord}}
\title{Effective heights of sums of elliptic abscissas}
\author{}\date{8 November 2026}
\begin{document}\maketitle
\begin{abstract}
For an elliptic curve $y^2=x^3-Dx$ over $\Q(\sqrt2)$ with
integral $D$, we bound the height lost when adding the
abscissas of finitely many integer multiples of a fixed point.
If the abscissas are positive at both real embeddings, the
loss is at most an explicit constant times the logarithm of
the largest index. The proof uses formal logarithms only as
identities: it requires no lower bound for a transcendental
linear form. Relative to any individual abscissa, for four
multiples the loss is at most $24\log B+\log864$.
Applied to the fixed elliptic source,
this gives a numerical bound in terms of the two rational
components of the sum, valid for all nonzero integer indices.
\end{abstract}

\section{Heights and the main statement}

Put $K=\Q(\sqrt2)$. At a finite place $v$ above a rational
prime $p$, let $\ord_v$ have value group $\Z$, let $f_v$ be
the residue degree, and put
$\lambda_v=f_v\log p/2$. Thus $\ord_v(p)=e_v$ and
$\sum_{v\mid p}e_vf_v=2$. For $a_1,\ldots,a_r\in K$ define
\begin{align*}
 h(1,a_1,\ldots,a_r)
 &=\sum_{v\nmid\infty}\lambda_v
          \max(0,-\ord_v(a_1),\ldots,-\ord_v(a_r))
 \\&\quad+\frac12\sum_{\sigma:K\hookrightarrow\mathbb R}
          \log\max(1,|\sigma(a_1)|,\ldots,|\sigma(a_r)|).
\end{align*}
For one entry this is the absolute logarithmic Weil height
$h(a)=h(1,a)$.

Let $D\in\Z[\sqrt2]\setminus\{0\}$ and
\[
 E_D:\quad y^2=x^3-Dx.
\]
Let $P\in E_D(K)$ have infinite order. For nonzero integers
$n_1,\ldots,n_r$ set
\[
 x_i=x([n_i]P),\qquad Z=\sum_{i=1}^r x_i,
 \qquad B=\max_i|n_i|\ge1.
\]
The points $[n_i]P$ are neither the origin nor a point of
order two, so every $x_i$ is defined and nonzero.

\begin{theorem}\label{thm:main}
Assume $\sigma(x_i)>0$ at both real embeddings of $K$.
Put
\[
 A_r=(r-1)2^{r-1},\qquad
 B_r=\sum_{k=1}^r\binom{r-1}{k-1}\log k.
\]
Then
\begin{equation}\label{eq:individual-height}
 h(Z)\ge\max_i h(x_i)-A_r\log B-B_r.
\end{equation}
In particular, for $r=4$,
\begin{equation}\label{eq:four}
 h(Z)\ge \max_i h(x_i)-24\log B-\log864.
\end{equation}
There is also a comparison with the joint coordinate height.
Put
\[
 C_r=2r+2\sum_{k=1}^r\binom rk(k-1)
     =(r-2)2^r+2r+2,
 \qquad
 D_r=\sum_{k=1}^r\binom rk\log k.
\]
\begin{equation}\label{eq:height}
 h(1,x_1,\ldots,x_r)-C_r\log B-D_r
 \ \le\ h(Z)\ \le\ h(1,x_1,\ldots,x_r)+\log r.
\end{equation}
All constants in these bounds are independent of $D$, $P$,
and the indices, subject to the stated hypotheses.
\end{theorem}

The positivity hypothesis is automatic when $P\in2E_D(K)$.
Indeed the duplication formula is
\begin{equation}\label{eq:duplicate}
 x(2Q)=\left(\frac{x(Q)^2+D}{2y(Q)}\right)^2.
\end{equation}
For a point $Q$ of infinite order the displayed denominator
and numerator are nonzero: otherwise $2Q$ would be the origin
or $(0,0)$. Hence its image under each real embedding is a
strictly positive square.

\section{An integral local principal part}

We give the local calculation at every finite place, including
places of bad reduction. Put $t=-x/y$, the parameter at the
origin. The integral Weierstrass equation defines its usual
formal group over each valuation ring. Its points with
$t$ in the maximal ideal form an open subgroup $E_1(K_v)$;
no good-reduction hypothesis is needed at the smooth origin.
For a finite nonzero point, membership in $E_1(K_v)$ is
equivalent to $\ord_v(x)<0$. In that case
\begin{equation}\label{eq:valuations}
 \ord_v(x)=-2\ord_v(t),\qquad
 \ord_v(y)=-3\ord_v(t).
\end{equation}
For example, these equalities follow directly by comparing
the terms of $y^2=x^3-Dx$ when $\ord_v(x)<0$.

\begin{lemma}\label{lem:local}
Let $v$ be any finite place of $K$. The formal logarithm
$\ell_v$ on $E_1(K_v)$ satisfies, for $a=\ord_v(t)>0$,
\[
 \ord_v(\ell_v(t))=a,
 \qquad x(t)-\ell_v(t)^{-2}\in\mathcal O_v.
\]
It is injective on $E_1(K_v)$.
\end{lemma}
\begin{proof}
The equation in the parameter $t$ gives
\begin{equation}\label{eq:x-series}
 x(t)=\frac{1+\sqrt{1+4Dt^4}}{2t^2}
 =t^{-2}+Dt^2-D^2t^6+2D^3t^{10}-\cdots.
\end{equation}
The coefficients after $t^{-2}$ are signed Catalan numbers
times powers of $D$, and are integral. Thus
$x(t)-t^{-2}\in\mathcal O_v$.

The invariant differential and its integral are
\begin{align}
 \frac{dx}{2y}&=(1+4Dt^4)^{-1/2}\,dt,\label{eq:differential}\\
 \ell_v(t)&=\sum_{k\ge0}
  \frac{(-1)^k\binom{2k}{k}D^k}{4k+1}\,t^{4k+1}.
                                                        \label{eq:log}
\end{align}
If $v\mid2$, every denominator $4k+1$ is a unit; the terms
with $k\ge1$ have valuation at least $5a$.
If $v\nmid2$, the extension $K/\Q$ is unramified at $v$.
For $k\ge1$ and odd $p$ one has
\[
 v_p(4k+1)\le4k-2,
\]
because $3^{4k-2}\ge4k+1$ (by induction, starting with
$9\ge5$). The $k$-th term of~\eqref{eq:log} consequently
has valuation at least
\[
 (4k+1)a-v_p(4k+1)\ge3a.
\]
Moreover $v_p(4k+1)\le\log_p(4k+1)=o(k)$, so the valuations
of the terms tend to infinity. At places above two this
also follows immediately from integrality of the coefficients.
Hence at every finite place
$\ord_v(\ell_v(t)-t)\ge3a$, and in particular
$\ord_v(\ell_v(t))=a$. Writing $\ell_v(t)=t(1+\eta)$
with $\ord_v(\eta)\ge2a$ gives
\[
 \ell_v(t)^{-2}-t^{-2}
 =t^{-2}\frac{-2\eta-\eta^2}{(1+\eta)^2}\in\mathcal O_v.
\]
Together with~\eqref{eq:x-series}, this proves the claimed
principal part. The integral of the invariant differential
is a formal-group homomorphism. The valuation equality
shows that its only zero in $E_1(K_v)$ is the origin, which
proves injectivity.
\end{proof}

Since $E_D(K_v)$ is compact and $E_1(K_v)$ is an open subgroup,
their quotient is finite. Choose $m>0$ with $[m]P\in E_1(K_v)$
and define
\[
 L_v=\frac{\ell_v([m]P)}m.
\]
The formal-group homomorphism shows this is independent of
the choice of $m$. It is nonzero, since $P$ has infinite order
and Lemma~\ref{lem:local} is injective. Whenever
$[n]P\in E_1(K_v)$,
\begin{equation}\label{eq:rank-one}
 \ell_v([n]P)=nL_v,
 \qquad x([n]P)=\frac1{n^2L_v^2}+\varepsilon_n,
 \quad \varepsilon_n\in\mathcal O_v.
\end{equation}
This identity is the only use of a local logarithm in the
height argument.

\section{Summing the local losses}

For each nonempty subset $I\subseteq\{1,\ldots,r\}$ put
\[
 T_I=\sum_{i\in I}\frac1{n_i^2}\in\Q_{>0}.
\]
Write $T_I=A_I/B_I$ in lowest terms, with positive integers
$A_I,B_I$.

\begin{lemma}\label{lem:fixed-target}
Fix an index $i$. At every finite place where $[n_i]P$ is
formal, let $J$ be the set of all formal indices and put
\[
 C_J=\sum_{j\in J}\prod_{k\in J\setminus\{j\}}n_k^2
       \in\Z_{>0}.
\]
The loss of the local height of $x_i$ upon passing to $Z$ is
at most $\ord_v(C_J)$.
\end{lemma}
\begin{proof}
Write $M_i=-\ord_v(x_i)>0$. The principal-part identity is
\[
 Z=(n_iL_v)^{-2}A_{J,i}+\varepsilon,
 \qquad
 A_{J,i}=n_i^2T_J
       =\frac{C_J}{\prod_{j\in J\setminus\{i\}}n_j^2},
 \qquad\varepsilon\in\mathcal O_v.
\]
If $\ord_v(A_{J,i})<M_i$, the local height of $Z$ is
$M_i-\ord_v(A_{J,i})$; otherwise it is zero. The loss is
therefore at most $\max(0,\ord_v(A_{J,i}))$.
All the indices are integers, so
$\ord_v(A_{J,i})\le\ord_v(C_J)$; the latter is nonnegative.
\end{proof}

\begin{lemma}\label{lem:loss}
At a finite place $v$, let
\[
 M_v=\max(0,-\ord_v(x_1),\ldots,-\ord_v(x_r)),\quad
 m_v=\max(0,-\ord_v(Z)).
\]
If $M_v>0$, let $I_v$ be the nonempty set of indices whose
points belong to $E_1(K_v)$, and choose $j\in I_v$ with
$\ord_v(n_j)=\max_{i\in I_v}\ord_v(n_i)$. Then
\[
 0\le M_v-m_v\le
 \ord_v(n_j^2T_{I_v})
 \le2\max_i\ord_v(n_i)+\max(0,\ord_v(T_{I_v})).
\]
\end{lemma}
\begin{proof}
The remaining abscissas are integral at $v$. By
\eqref{eq:rank-one},
\[
 Z=\frac{T_{I_v}}{L_v^2}+\varepsilon,
 \qquad\varepsilon\in\mathcal O_v,
 \qquad M_v=2\ord_v(n_jL_v)>0.
\]
The number
$c=n_j^2T_{I_v}=\sum_{i\in I_v}(n_j/n_i)^2$
is nonzero and integral at $v$. If $\ord_v(c)<M_v$, the
principal part has negative valuation and
$m_v=M_v-\ord_v(c)$. Otherwise the principal part and $Z$
are integral, and $m_v=0$. Thus the loss is
$\min(M_v,\ord_v(c))$, proving the inequalities.
\end{proof}

\begin{proof}[Proof of Theorem~\ref{thm:main}]
Fix an index $i$. Where $[n_i]P$ is not formal, its finite
local height is zero and there is no positive loss. At every
other place Lemma~\ref{lem:fixed-target} applies. As $J$
ranges over the subsets containing $i$, the sum of its
contributions is at most
\[
 \sum_{J\ni i}\sum_{v\nmid\infty}
                  \lambda_v\ord_v(C_J)
 =\sum_{J\ni i}\log C_J.
\]
For $|J|=k$ one has $C_J\le kB^{2(k-1)}$, and there are
$\binom{r-1}{k-1}$ such subsets. Consequently this total
is at most
\[
 2\sum_{k=1}^r\binom{r-1}{k-1}(k-1)\log B+B_r
 =(r-1)2^{r-1}\log B+B_r.
\]
At a real embedding the positive sum is at least its
$i$-th summand, so there is no archimedean loss.
This proves~\eqref{eq:individual-height} for every $i$.
The values $A_4=24$ and
$B_4=3\log2+3\log3+\log4=\log864$ give~\eqref{eq:four}.

The ultrametric inequality shows $m_v\le M_v$. At places
with $M_v=0$ there is no loss. Sum Lemma~\ref{lem:loss},
using its last inequality, to obtain
\begin{align*}
 \sum_{v\nmid\infty}\lambda_v(M_v-m_v)
 &\le2\sum_{i=1}^r\log|n_i|
    +\sum_{\varnothing\ne I\subseteq\{1,\ldots,r\}}
       \log A_I.                                      \tag{*}
\end{align*}
Here, for any positive rational $a=A/B$ in lowest terms,
\[
 \sum_{v\nmid\infty}\lambda_v\max(0,\ord_v(a))=
 \sum_p\log p\,\max(0,v_p(a))=\log A.
\]
Also $\max_i\ord_v(n_i)\le\sum_i\ord_v(n_i)$ because
all the $n_i$ are integers. This proves~$(*)$, even though
the subset $I_v$ varies with the place: each of the finitely
many subsets contributes at most its full positive-valuation
sum.

If $|I|=k$, a common denominator for $T_I$ gives
\[
 A_I\le\sum_{i\in I}\prod_{j\in I\setminus\{i\}}n_j^2
       \le kB^{2(k-1)}.
\]
Together with $\log|n_i|\le\log B$, this bounds~$(*)$
by $C_r\log B+D_r$. At each real embedding all $x_i$ are
positive, so
\[
 \max(1,\sigma(x_1),\ldots,\sigma(x_r))
 \le\max(1,\sigma(Z))
 \le r\max(1,\sigma(x_1),\ldots,\sigma(x_r)).
\]
Combining the finite and infinite places proves
\eqref{eq:height}. Its four-coordinate constants are
$C_4=42$ and $D_4=\log20736$.
\end{proof}

\section{The fixed arithmetic source}

Take
\[
 D=11-6\sqrt2=(3-\sqrt2)^2,
 \qquad R_0=(1+\sqrt2,2\sqrt2),
\]
\[
 P=2R_0=
 \left(\frac{57}8-\frac72\sqrt2,
       -\frac{165}8+\frac{417}{32}\sqrt2\right).
\]
The equation and the doubling identity are checked directly.
At the place $\pi=\sqrt2$ one has
$\ord_\pi(x(P))=-6$ and $\ord_\pi(y(P))=-9$; hence
$\ord_\pi(t(P))=3$.
Lemma~\ref{lem:local} proves that $\ell_\pi(P)\ne0$, so
$P$ has infinite order. Formula~\eqref{eq:duplicate} proves
the positivity needed in Theorem~\ref{thm:main} for every
nonzero integer index.

The arithmetic source uses the index-height assertion
\begin{equation}\label{eq:source-height}
                 h(x([n]P))\ge n^2-6\qquad(n\in\Z\setminus\{0\}).
\end{equation}
We state the following application explicitly under that
retained assertion; its precise source identifier is
\nolinkurl{eq:o8-indexheight}, as documented in~\cite{Arithmetic}.
The sum-height theorem above is independent of this assertion.

\begin{corollary}\label{cor:numerical}
Assume~\eqref{eq:source-height}. For nonzero integers
$n_1,n_2,n_3,n_4$, put
\[
 Z=\sum_{i=1}^4x([n_i]P)=S_0+S_1\sqrt2,
 \qquad S_0,S_1\in\Q,
 \qquad B=\max_i|n_i|.
\]
Then
\begin{equation}\label{eq:bounds}
 B^2\le2h(Z)+80
       \le2h(S_0)+2h(S_1)+83.
\end{equation}
No restriction on the parity of the indices is imposed.
\end{corollary}
\begin{proof}
Theorem~\ref{thm:main} gives
\[
 h(Z)\ge B^2-24\log B-6-\log864.
\]
For $B>0$, the inequality $\log t\le t-1$ with
$t=B^2/24$ gives
\[
 24\log B-\frac{B^2}2
 \le12\log24-12<27.
\]
Since $\log864<7$, the total subtracted constant is
less than $40$. Thus $h(Z)\ge B^2/2-40$.
Finally the elementary height inequalities for addition and
multiplication give
\[
 h(S_0+S_1\sqrt2)
 \le h(S_0)+h(S_1)+\tfrac32\log2.
\]
Because $3\log2<3$, the second inequality follows after
increasing the additive constant from $80$ to $83$.
The statements $\log24<13/4$ and $\log864<7$ can, if
desired, be verified by finite positive Taylor sums for the
exponential; exact certificates are supplied with this paper.
\end{proof}

The two rational components $S_0,S_1$ are linear combinations
of the existing abscissa components. Thus this aggregate
requires no multiplication coordinates. Inequality~\eqref{eq:bounds}
is the height input for a construction retaining the three
scales $1$, $S_0^2$, and $(1+S_1^2)^2$.
It does not assert an effective bound in terms of $h(S_0)$
alone; taking the field trace can introduce additional
cancellation between conjugate blocks.

\section{A number-field extension}

\begin{corollary}
Let $F$ be a totally real number field unramified at every
odd rational prime. Let $D\in\mathcal O_F\setminus\{0\}$
and let $P\in2E_D(F)$ have infinite order. Then both
inequalities~\eqref{eq:individual-height} and~\eqref{eq:height}
hold for the absolute heights over $F$, with the same
constants $A_r,B_r,C_r,D_r$.
\end{corollary}
\begin{proof}
At every odd place the valuation estimate in
Lemma~\ref{lem:local} uses only that the ramification index
is one. At places above two, all coefficients of the formal
logarithm are integral, regardless of the ramification
index. The lemma therefore holds over $F$ at every finite
place. In the product-formula calculation replace
$f_v\log p/2$ by $f_v\log p/[F:\Q]$; the contribution of
a rational positive integer remains its usual logarithm.
At real embeddings positivity follows from doubling as
before. The two proofs are otherwise unchanged.
\end{proof}

\section{Verification data}

The companion program
\nolinkurl{check_effective_elliptic_sums.py} uses exact rational
arithmetic. It verifies the curve and doubling identities,
the stated valuation of $P$, the first twelve coefficients
of the formal identities, the combinatorial constants, and
finite rational certificates for the two exponential
inequalities used above. The local estimates and the
summation over all places are proved in the text, rather
than inferred from a finite computation.

\begin{thebibliography}{9}
\bibitem{Silverman}
J.~H.~Silverman, \emph{The Arithmetic of Elliptic Curves},
second edition, Graduate Texts in Mathematics 106,
Springer, 2009. Chapters~IV and~VII supply the standard
formal-group and local-group terminology used here.
\url{https://doi.org/10.1007/978-0-387-09494-6}.
\bibitem{Arithmetic}
\emph{Arithmetic dependencies of the optimized rational
constructions}, 8 November 2026, Section~2,
\nolinkurl{arithmetic_dependencies_20261108_v8.tex},
project replication archive.
\end{thebibliography}
\end{document}
